\section{团队分工}

本实验由四名组员协作完成，整体工作包括数据预处理、语音信号分析、特征提取、传统机器学习模型、深度学习模型、噪声鲁棒性测试、结果可视化、报告撰写和展示材料整理。各成员在完成各自模块的基础上，也对整体实验流程、结果分析和报告修改进行了讨论与协作。

\begin{table}[H]
  \centering
  \caption{团队成员分工}
  \label{tab:team-contribution}
  \small
  \begin{tabular}{p{0.16\linewidth}p{0.76\linewidth}}
    \toprule
    成员 & 主要工作 \\
    \midrule
    于沛龙 &
    负责传统方法基线流程搭建，包括语音预处理、端点检测、MFCC 特征提取、KNN 基线模型实现与测试；参与噪声鲁棒性实验、结果图表整理、PPT 制作与实验报告统稿。 \\
    \addlinespace

    景序 &
    负责传统方法优化部分，设计并实现基于 MFCC + $\Delta$ + $\Delta^2$ 帧级特征的 GMM 序列模型；完成 GMM 纯净语音测试、噪声鲁棒性重复实验、准确率--SNR 曲线绘制及相关结果分析。 \\
    \addlinespace

    封宇凡 &
    负责深度学习方法部分，设计并训练 CNN 模型；完成 MFCC 动态特征图输入构建、网络结构设计、模型训练参数调优、纯净语音测试和混合噪声细粒度 SNR 测试。 \\
    \addlinespace

    孟德轩 &
    参与数据处理、实验结果核对、图表整理和报告材料撰写；协助完成实验流程梳理、不同方法结果对照、报告格式检查以及展示内容整理。 \\
    \bottomrule
  \end{tabular}
\end{table}

在协作过程中，小组首先确定整体技术路线：
以 MFCC 为核心声学特征，分别构建 KNN、GMM 和 CNN 三种识别方法，
并统一使用 speaker1--5 作为训练集、speaker6--7 作为验证集。
随后，各成员分别完成对应模型的实现与测试，并将输出结果统一整理为混淆矩阵、准确率表格和准确率--SNR 曲线。
最后，小组共同对实验结果进行比较分析，讨论不同方法在小样本条件和混合噪声条件下的优劣，完成最终实验报告和课程展示材料。
